\ifdefined\gkver\else\def\gkver{student}\fi
\documentclass[\gkver]{gaokaozhenti}
\begin{document}

\chapter{2004年广东广西}

\begin{enumerate}
\item
%% number: 1
%% paperName: 2004年广东广西高考第 1 题 0 分
%% typeId: 1
%% type: 单选题
%% chapter: 原子和原子核
%% point: 玻尔模型
%% method: 无
%% score: 0
%% degree: 650
%% duplicateId: 0
%% body:
图示为氢原子的能级图，用光子能量为 $13.07\UeV$ 的光照射一群处于基态的氢原子，可能观测到氢原子发射的不同波长有多少种？\xzanswer{B}
\onepicture[w=5.3cm]{figs/01.png}
\fourchoices[answer=B]
{15}
{10}
{4}
{1}
%% answer: B
%% memo:
\memoanswer{
由题图知，$13.07\UeV$ 恰好为 $n=5$ 与 $n=1$ 的能级差，从而使基态氢原子跃迁到 $n=5$ 的激发态，这些氢原子再向较低能级跃迁时，可产生不同波长的光共 $C_{5}^{2}=10$ 种，故 B 正确。
}

\item
%% number: 2
%% paperName: 2004年广东广西高考第 2 题 4 分
%% typeId: 2
%% type: 多选题
%% chapter: 气体性质
%% point: 分子力
%% method: 无
%% score: 4
%% degree: 650
%% duplicateId: 0
%% body:
下列说法哪些是正确的\xzanswer{AD}
\fourchoices[answer=AD]
{水的体积很难被压缩，这是分子间存在斥力的宏观表现}
{气体总是很容易充满容器，这是分子间存在斥力的宏观表现}
{两个相同的半球壳吻合接触，中间抽成真空（马德堡半球），用力很难拉开，这是分子间存在吸引力的宏观表现}
{用力拉铁棒的两端，铁棒没有断，这是分子间存在吸引力的宏观表现}
%% answer: AD
%% memo:
\memoanswer{
液体体积很难压缩，说明分子间存在斥力；固体很难被拉断，说明分子间存在引力，故 A、D 正确。气体容易充满容器是分子热运动的结果，抽成真空的马德堡半球很难分开是大气压强作用的结果，故 B、C 错误。
}

\item
%% number: 3
%% paperName: 2004年广东广西高考第 3 题 4 分
%% typeId: 2
%% type: 多选题
%% chapter: 机械波
%% point: 波的多解性
%% method: 无
%% score: 4
%% degree: 650
%% duplicateId: 0
%% body:
一列简谐波沿一直线向左运动，当直线上某质点 a 向上运动到达最大位移时，a 点右方相距 $0.15\Um$ 的 b 点刚好向下运动到最大位移处，则这列波的波长可能是\xzanswer{BD}
\onepicture[w=8.0cm]{\includesvg{figs/03}}
\fourchoices[answer=BD]
{$0.6\Um$}
{$0.3\Um$}
{$0.2\Um$}
{$0.1\Um$}
%% answer: BD
%% memo:
\memoanswer{
由题意知，$0.15=\left(n+\frac{1}{2}\right)\lambda$（$n=0,1,2,\dots$），得 $\lambda=\frac{0.3}{2n+1}\Um$（$n=0,1,2,\dots$），当 $n=0$ 时，$\lambda=0.3\Um$；当 $n=1$ 时，$\lambda=0.1\Um$，故 B、D 正确。
}

\item
%% number: 4
%% paperName: 2004年广东广西高考第 4 题 4 分
%% typeId: 1
%% type: 单选题
%% chapter: 气体性质
%% point: 热力学第一定律
%% method: 无
%% score: 4
%% degree: 650
%% duplicateId: 0
%% body:
下列说法正确的是\xzanswer{D}
\fourchoices[answer=D]
{机械能全部变成内能是不可能的}
{第二类永动机不可能制造成功的原因是因为能量既不会凭空产生，也不会凭空消失，只能从一个物体转移到另一个物体，或从一种形式转化成另一种形式}
{根据热力学第二定律可知，热量不可能从低温物体传到高温物体}
{从单一热源吸收的热量全部变成功是可能的}
%% answer: D
%% memo:
\memoanswer{
（广西卷）由热力学第一定律 $\Delta U=W+Q$ 可知，$\Delta U$ 由 $W+Q$ 共同决定，故 A、C 错误，D 正确；机械能可以通过克服摩擦力做功全部转化为内能，故 B 错误。

（广东卷）机械能可通过克服摩擦力做功全部转化为内能，故 A 错误；第二类永动机没有违背能量守恒定律，而是违反了热力学第二定律，因而不能制成，故 B 错误；由热力学第二定律可知，热量不能自发地从低温传到高温物体，但在外部条件的影响下，能从低温传到高温物体，故 C 错误；由热力学第二定律可知，从单一热源吸收热量全部变成功也是可能的，但要发生其他变化，故 D 正确。
}

\item
%% number: 5
%% paperName: 2004年广东广西高考第 5 题 4 分
%% typeId: 1
%% type: 单选题
%% chapter: 原子和原子核
%% point: 质能方程
%% method: 无
%% score: 4
%% degree: 650
%% duplicateId: 0
%% body:
中子 n、质子 p、氘核 D 的质量分别为 $m_{n}$、$m_{p}$、$m_{D}$。现用光子能量为 $E$ 的 $\gamma$ 射线照射静止氘核使之分解，反应的方程为 $\gamma+\mathrm{D}=\mathrm{p}+\mathrm{n}$，若分解后中子、质子的动能可视为相等，则中子的动能是\xzanswer{C}
\fourchoices[answer=C]
{$\frac{1}{2}[(m_{D}-m_{p}-m_{n})c^{2}-E]$}
{$\frac{1}{2}[(m_{D}+m_{p}-m_{n})c^{2}+E]$}
{$\frac{1}{2}[(m_{D}-m_{p}-m_{n})c^{2}+E]$}
{$\frac{1}{2}[(m_{D}+m_{p}-m_{n})c^{2}-E]$}
%% answer: C
%% memo:
\memoanswer{
由能量守恒得 $E+\Delta mc^{2}=2E_{k}$，$\Delta m=m_{D}-m_{p}-m_{n}$，得 $E_{k}=\frac{1}{2}[E+(m_{D}-m_{p}-m_{n})c^{2}]$，故 C 正确。
}

\item
%% number: 6
%% paperName: 2004年广东广西高考第 6 题 4 分
%% typeId: 1
%% type: 单选题
%% chapter: 光学
%% point: 光电效应
%% method: 无
%% score: 4
%% degree: 650
%% duplicateId: 0
%% body:
分别用波长为 $\lambda$ 和 $\frac{3}{4}\lambda$ 的单色光照射同一金属板，发出的光电子的最大初动能之比为 $1:2$，以 $h$ 表示普朗克常量，$c$ 表示真空中的光速，则此金属板的逸出功为\xzanswer{B}
\fourchoices[answer=B]
{$\frac{hc}{2\lambda}$}
{$\frac{2hc}{3\lambda}$}
{$\frac{3}{4}hc\lambda$}
{$\frac{4h\lambda}{5c}$}
%% answer: B
%% memo:
\memoanswer{
由光电效应方程 $E_{k}=h\nu-W=\frac{hc}{\lambda}-W$，对波长 $\lambda$ 有 $E_{k1}=\frac{hc}{\lambda}-W$，对波长 $\frac{3}{4}\lambda$ 有 $E_{k2}=\frac{hc}{\frac{3}{4}\lambda}-W=\frac{4hc}{3\lambda}-W$。由 $E_{k1}:E_{k2}=1:2$ 得 $E_{k2}=2E_{k1}$，即 $\frac{4hc}{3\lambda}-W=2\left(\frac{hc}{\lambda}-W\right)$，解得 $W=\frac{2hc}{3\lambda}$，故 B 正确。
}

\item
%% number: 7
%% paperName: 2004年广东广西高考第 7 题 4 分
%% typeId: 1
%% type: 单选题
%% chapter: 力物体的平衡
%% point: 三力平衡
%% method: 无
%% score: 4
%% degree: 650
%% duplicateId: 0
%% body:
用三根轻绳将质量为 $m$ 的物块悬挂在空中，如图所示。已知 ac 和 bc 与竖直方向的夹角分别为 $30^{\circ}$ 和 $60^{\circ}$，则 ac 绳和 bc 绳中的拉力分别为\xzanswer{A}
\onepicture[w=4.2cm]{\includesvg{figs/07}}
\fourchoices[answer=A]
{$\frac{\sqrt{3}}{2}mg$，$\frac{1}{2}mg$}
{$\frac{1}{2}mg$，$\frac{\sqrt{3}}{2}mg$}
{$\frac{\sqrt{3}}{4}mg$，$\frac{1}{2}mg$}
{$\frac{1}{2}mg$，$\frac{\sqrt{3}}{4}mg$}
%% answer: A
%% memo:
\memoanswer{
对结点 c 受力分析如图所示，由平衡条件得 $T_{ac}=mg\cos30^{\circ}=\frac{\sqrt{3}}{2}mg$，$T_{bc}=mg\sin30^{\circ}=\frac{1}{2}mg$，故 A 正确。

\onepicture[w=3.5cm]{figs/07_an.jpg}
}

\item
%% number: 8
%% paperName: 2004年广东广西高考第 8 题 4 分
%% typeId: 1
%% type: 单选题
%% chapter: 气体性质
%% point: 热力学第一定律
%% method: 无
%% score: 4
%% degree: 450
%% duplicateId: 0
%% body:
如图所示，密闭绝热的具有一定质量的活塞，活塞的上部封闭着气体，下部为真空，活塞与器壁的摩擦忽略不计，置于真空中的轻弹簧的一端固定于容器的底部。另一端固定在活塞上，弹簧被压缩后用绳扎紧，此时弹簧的弹性势能为 $E_{p}$（弹簧处于自然长度时的弹性势能为零），现绳突然断开，弹簧推动活塞向上运动，经过多次往复运动后活塞静止，气体达到平衡态，经过此过程\xzanswer{D}
\onepicture[w=4.0cm]{\includesvg{figs/08}}
\fourchoices[answer=D]
{$E_{p}$ 全部转换为气体的内能}
{$E_{p}$ 一部分转换成活塞的重力势能，其余部分仍为弹簧的弹性势能}
{$E_{p}$ 全部转换成活塞的重力势能和气体的内能}
{$E_{p}$ 一部分转换成活塞的重力势能，一部分转换为气体的内能，其余部分仍为弹簧的弹性势能}
%% answer: D
%% memo:
\memoanswer{
由题意，断开绳子，活塞最终静止后的位置高于初始位置，$E_{p}$ 的能量转化有三种形式：活塞的重力势能、气体的内能和弹簧的弹性势能，故 D 正确。
}

\item
%% number: 9
%% paperName: 2004年广东广西高考第 9 题 4 分
%% typeId: 1
%% type: 单选题
%% chapter: 直线运动
%% point: 竖直上下抛运动
%% method: 无
%% score: 4
%% degree: 650
%% duplicateId: 0
%% body:
一杂技演员，用一只手抛球。他每隔 $0.40\Us$ 抛出一球，接到球便立即把球抛出，已知除抛、接球的时刻外，空中总有四个球，将球的运动看作是竖直方向的运动，球到达的最大高度是（高度从抛球点算起）\xzanswer{C}
\fourchoices[answer=C]
{$1.6\Um$}
{$2.4\Um$}
{$3.2\Um$}
{$4.0\Um$}
%% answer: C
%% memo:
\memoanswer{
如图所示，当某一球刚上抛时，最高点有一球，则 $h=\frac{1}{2}gt^{2}=\frac{1}{2}\times10\times(2\times0.4)^{2}\Um=3.2\Um$，故 C 正确。

\onepicture[h=4.0cm]{figs/09_an.jpg}
}

\item
%% number: 10
%% paperName: 2004年广东广西高考第 10 题 4 分
%% typeId: 2
%% type: 多选题
%% chapter: 电场
%% point: 带电粒子的直线运动
%% method: 无
%% score: 4
%% degree: 450
%% duplicateId: 0
%% body:
在场强为 $E$ 的匀强电场中固定放置两个小球 1 和 2，它们的质量相等，电荷分别为 $q_{1}$ 和 $-q_{2}$（$q_{1}\neq q_{2}$）。球 1 和球 2 的连线平行于电场线，如图。现同时放开 1 球和 2 球，于是它们开始在电场力的作用下运动，如果球 1 和球 2 之间的距离可以取任意有限值，则两球刚被放开时，它们的加速度可能是\xzanswer{ABC}
\onepicture[w=4.8cm]{\includesvg{figs/10}}
\fourchoices[answer=ABC]
{大小相等，方向相同}
{大小不等，方向相反}
{大小不等，方向相同}
{大小相等，方向相反}
%% answer: ABC
%% memo:
\memoanswer{
设两球距离为 $r$，取向右为正，由牛顿第二定律，对球 1 有 $q_{1}E-k\frac{q_{1}q_{2}}{r^{2}}=ma_{1}$，对球 2 有 $k\frac{q_{1}q_{2}}{r^{2}}-q_{2}E=ma_{2}$，由于 $q_{1}\neq q_{2}$ 且 $r$ 为任意有限值，可知 A、B、C 正确。要使 D 成立，必须 $q_{1}=q_{2}$，与题意矛盾，故 D 错误。
}

\item
%% number: 11
%% paperName: 2004年广东广西高考第 11 题 8 分
%% typeId: 4
%% type: 实验题
%% chapter: 光学
%% point: 测定玻璃折射率
%% method: 无
%% score: 8
%% degree: 650
%% duplicateId: 0
%% body:
如图，画有直角坐标系 $Oxy$ 的白纸位于水平桌面上，M 是放在白纸上的半圆形玻璃砖，其底面的圆心在坐标的原点，直边与 $x$ 轴重合，OA 是画在纸上的直线，$P_{1}$、$P_{2}$ 为竖直地插在直线 OA 上的两枚大头针，$P_{3}$ 是竖直地插在纸上的第三枚大头针，$\alpha$ 是直线 OA 与 $y$ 轴正方向的夹角，$\beta$ 是直线 $OP_{3}$ 与 $y$ 轴负方向的夹角，只要直线 OA 画得合适，且 $P_{3}$ 的位置取得正确，测得角 $\alpha$ 和 $\beta$，便可求得玻璃的折射率。

某学生在用上述方法测量玻璃的折射率，在他画出的直线 OA 上竖直插上了 $P_{1}$、$P_{2}$ 两枚大头针，但在 $y<0$ 的区域内，不管眼睛放在何处，都无法透过玻璃砖看到 $P_{1}$、$P_{2}$ 的像，他应该采取的措施是 \tkanswer[8cm]{在白纸上另画一条与 $y$ 轴正方向的夹角较小的直线 OA，把大头针 $P_{1}$、$P_{2}$ 竖直地插在所画的直线上，直到在 $y<0$ 的区域内透过玻璃砖能看到 $P_{1}$、$P_{2}$ 的像}。若他已透过玻璃砖看到了 $P_{1}$、$P_{2}$ 的像，确定 $P_{3}$ 位置的方法是 \tkanswer[5cm]{插上 $P_{3}$ 后，$P_{3}$ 刚好能挡住 $P_{1}$、$P_{2}$ 的像}。若他已正确地测得了的 $\alpha$、$\beta$ 的值，则玻璃的折射率 $n=$ \tkanswer[3cm]{$\frac{\sin\beta}{\sin\alpha}$}。

\onepicture[align=right, w=5.3cm]{\includesvg{figs/11}}
%% answer: 
%% memo:
\memoanswer{
在白纸上另画一条与 $y$ 轴正方向的夹角较小的直线 OA，把大头针 $P_{1}$、$P_{2}$ 竖直地插在所画的直线上，直到在 $y<0$ 的区域内透过玻璃砖能看到 $P_{1}$、$P_{2}$ 的像。

插上 $P_{3}$ 后，$P_{3}$ 刚好能挡住 $P_{1}$、$P_{2}$ 的像。

$n=\frac{\sin\beta}{\sin\alpha}$。
}

\item
%% number: 12
%% paperName: 2004年广东广西高考第 12 题 12 分
%% typeId: 4
%% type: 实验题
%% chapter: 电路
%% point: 测电动势与内阻
%% method: 无
%% score: 12
%% degree: 550
%% duplicateId: 0
%% body:
图中 $R$ 为已知电阻，$R_{x}$ 为待测电阻，K$_{1}$ 为单刀单掷开关，K$_{2}$ 为单刀双掷开关，V 为电压表（内阻极大），$E$ 为电源（电阻不可忽略）。现用图中电路测量电源电动势 $\varepsilon$ 及电阻 $R_{x}$。

\begin{enumerate}
	\item
	写出操作步骤：\tkanswer[6cm]{}
	\item
	由 $R$ 及测得的量，可测得 $\varepsilon=$ \tkanswer[1.5cm]{}，$R_{x}=$ \tkanswer[2.5cm]{}。
\end{enumerate}

\onepicture[align=right, w=6.6cm]{figs/12.png}
%% answer: 
\jdanswer{
\begin{enumerate}
	\item
	①K$_{1}$ 断开，K$_{2}$ 接到 a 端，记下电压表的读数 $U_{1}$；②K$_{2}$ 仍接到 a 端，闭合 K$_{1}$，记下电压表的读数 $U_{2}$；③K$_{1}$ 仍闭合，K$_{2}$ 接到 b 端，记下电压表的读数 $U_{3}$。
	\item
	$U_{1}$，$\frac{U_{3}}{U_{2}-U_{3}}R$。
\end{enumerate}
}
%% memo:
\memoanswer{
（1）①K$_{1}$ 断开，K$_{2}$ 接到 a 端，记下电压表的读数 $U_{1}$；②K$_{2}$ 仍接到 a 端，闭合 K$_{1}$，记下电压表的读数 $U_{2}$；③K$_{1}$ 仍闭合，K$_{2}$ 接到 b 端，记下电压表的读数 $U_{3}$。

（2）$U_{1}$，$\frac{U_{3}}{U_{2}-U_{3}}R$。
}

\item
%% number: 13
%% paperName: 2004年广东广西高考第 13 题 12 分
%% typeId: 6
%% type: 计算题
%% chapter: 电场
%% point: 库仑定律
%% method: 无
%% score: 12
%% degree: 650
%% duplicateId: 0
%% body:
已经证实，质子、中子都是由上夸克和下夸克的两种夸克组成的，上夸克带电为 $\frac{2}{3}e$，下夸克带电为 $-\frac{1}{3}e$，$e$ 为电子所带电量的大小，如果质子是由三个夸克组成的，且各个夸克之间的距离都为 $l$，$l=1.5\times10^{-15}\Um$，试计算质子内相邻两个夸克之间的静电力（库仑力）。
%% answer: 
\jdanswer{
上夸克与上夸克间的静电力为斥力，大小为 $46\UN$；上夸克与下夸克间的静电力为吸引力，大小为 $23\UN$。
}
%% memo:
\memoanswer{
质子带电为 $+e$，所以它是由 2 个上夸克和 1 个下夸克组成的，三个夸克位于等边三角形的三个顶点处，这时上夸克与上夸克之间的静电力应为
$F_{1}=k\frac{\frac{2}{3}e\times\frac{2}{3}e}{l^{2}}=\frac{4}{9}k\frac{e^{2}}{l^{2}}$ ①

代入数值，得 $F_{1}=46\UN$，为斥力 ②

上夸克与下夸克之间的静电力为
$F_{2}=k\frac{\frac{2}{3}e\times\frac{1}{3}e}{l^{2}}=\frac{2}{9}k\frac{e^{2}}{l^{2}}$ ③

代入数值，得 $F_{2}=23\UN$，为引力 ④

评分参考：①式 4 分，②式 2 分；③式 4 分，④式 2 分。
}

\item
%% number: 14
%% paperName: 2004年广东广西高考第 14 题 14 分
%% typeId: 6
%% type: 计算题
%% chapter: 运动和力
%% point: 动量定理
%% method: 无
%% score: 14
%% degree: 650
%% duplicateId: 0
%% body:
一质量为 $m$ 的小球，以初速度 $v_{0}$ 沿水平方向射出，恰好垂直地射到一倾角为 $30^{\circ}$ 的固定斜面上，并立即反方向弹回。已知反弹速度的大小是入射速度大小的 $\frac{3}{4}$，求在碰撞中斜面对小球的冲量大小。
%% answer: 
\jdanswer{
$I=\frac{7}{2}mv_{0}$
}
%% memo:
\memoanswer{
小球在碰撞斜面前做平抛运动。设刚要碰撞斜面时小球速度为 $v$。由题意，$v$ 的方向与竖直线的夹角为 $30^{\circ}$，且水平分量仍为 $v_{0}$，如右图。由此得
$v=2v_{0}$ ①

碰撞过程中，小球速度由 $v$ 变为反向 $\frac{3}{4}v$。碰撞时间极短，可不计重力的冲量，由动量定理，斜面对小球的冲量为
$I=m\left(\frac{3}{4}v\right)+mv$ ②

由①②得
$I=\frac{7}{2}mv_{0}$ ③

\onepicture[w=4.8cm]{figs/14a_an.jpg}
}

\item
%% number: 14
%% paperName: 2004年广东广西高考第 14 题 14 分
%% typeId: 6
%% type: 计算题
%% chapter: 直线运动
%% point: 匀变速直线运动
%% method: 无
%% score: 14
%% degree: 450
%% duplicateId: 0
%% body:
一路灯距地面的高度为 $h$，身高为 $l$ 的人以速度 $v$ 匀速行走，如图所示。

\begin{enumerate}
	\item
	试证明人的头顶的影子作匀速运动；
	\item
	求人影的长度随时间的变化率。
\end{enumerate}

\onepicture[align=right, w=6.0cm]{figs/14b.png}
%% answer: 
\jdanswer{
\begin{enumerate}
	\item
	$OM=\frac{hv}{h-l}t$，因 OM 与时间成正比，故人头顶的影子作匀速运动；
	\item
	$k=\frac{lv}{h-l}$。
\end{enumerate}
}
%% memo:
\memoanswer{
（1）设 $t=0$ 时刻，人位于路灯的正下方 O 处，在时刻 $t$，人走到 s 处，根据题意有
$OS=vt$ ①

过路灯 P 和人头顶的直线与地面的交点 M 为 $t$ 时刻人头顶影子的位置，如图所示，OM 为人头顶影子到 O 点的距离，由几何关系，有
$\frac{h}{OM}=\frac{l}{OM-OS}$ ②

解①②得 $OM=\frac{hv}{h-l}t$ ③

因 OM 与时间成正比，故人头顶的影子作匀速运动。

（2）由图可知，在时刻 $t$，人影的长度为 $SM$，由几何关系，有 $SM=OM-OS$ ④

由①③④式得
$SM=\frac{lv}{h-l}t$ ⑤

可见影长 $SM$ 与时间 $t$ 成正比，所以影长随时间的变化率 $k=\frac{lv}{h-l}$ ⑥

\onepicture[w=6.5cm]{figs/14b_an.jpg}
}

\item
%% number: 15
%% paperName: 2004年广东广西高考第 15 题 15 分
%% typeId: 6
%% type: 计算题
%% chapter: 电磁感应
%% point: 双杆问题
%% method: 无
%% score: 15
%% degree: 450
%% duplicateId: 0
%% body:
如图，在水平面上有两条平行导电导轨 MN、PQ，导轨间距离为 $l$，匀强磁场垂直于导轨所在的平面（纸面）向里，磁感应强度的大小为 $B$，两根金属杆 1、2 摆在导轨上，与导轨垂直，它们的质量和电阻分别为 $m_{1}$、$m_{2}$ 和 $R_{1}$、$R_{2}$，两杆与导轨接触良好，与导轨间的动摩擦因数为 $\mu$，已知：杆 1 被外力拖动，以恒定的速度 $v_{0}$ 沿导轨运动；达到稳定状态时，杆 2 也以恒定速度沿导轨运动，导轨的电阻可忽略，求此时杆 2 克服摩擦力做功的功率。

\onepicture[align=right, w=8.0cm]{\includesvg{figs/15}}
%% answer: 
\jdanswer{
$P=\mu m_{2}g\left[v_{0}-\frac{\mu m_{2}g}{B^{2}l^{2}}(R_{1}+R_{2})\right]$
}
%% memo:
\memoanswer{
解法一：设杆 2 的运动速度为 $v$，由于两杆运动时，两杆和导轨构成的回路的磁通量发生变化，产生感应电动势
$\varepsilon=Bl(v_{0}-v)$ ①

感应电流
$I=\frac{\varepsilon}{R_{1}+R_{2}}$ ②

杆 2 运动受到的安培力等于摩擦力 $BIl=\mu m_{2}g$ ③

导体杆 2 克服摩擦力做功的功率 $P=\mu m_{2}gv$ ④

解得 $P=\mu m_{2}g\left[v_{0}-\frac{\mu m_{2}g}{B^{2}l^{2}}(R_{1}+R_{2})\right]$ ⑤

解法二：以 $F$ 表示拖动杆 1 的外力，$I$ 表示回路的电流，达到稳定时，对杆 1 有
$F-\mu m_{1}g-BIl=0$ ①

对杆 2 有 $BIl-\mu m_{2}g=0$ ②

外力的功率 $P_{F}=Fv_{0}$ ③

以 $P$ 表示杆 2 克服摩擦力做功的功率，则有 $P=P_{F}-I^{2}(R_{1}+R_{2})-\mu m_{1}gv_{0}$ ④

解得 $P=\mu m_{2}g\left[v_{0}-\frac{\mu m_{2}g}{B^{2}l^{2}}(R_{1}+R_{2})\right]$ ⑤

①②③④⑤各 3 分。
}

\item
%% number: 16
%% paperName: 2004年广东广西高考第 16 题 16 分
%% typeId: 6
%% type: 计算题
%% chapter: 曲线运动
%% point: 人造地球卫星
%% method: 无
%% score: 16
%% degree: 450
%% duplicateId: 0
%% body:
某颗地球同步卫星正下方的地球表面上有一观察者，他用天文望远镜观察被太阳光照射的此卫星，试问，春分那天（太阳光直射赤道）在日落 12 小时内有多长时间该观察者看不见此卫星？已知地球半径为 $R$，地球表面处的重力加速度为 $g$，地球自转周期为 $T$，不考虑大气对光的折射。
%% answer: 
\jdanswer{
$t=\frac{T}{\pi}\arcsin\left(\frac{4\pi^{2}R}{gT^{2}}\right)^{\frac{1}{3}}$
}
%% memo:
\memoanswer{
设所求的时间为 $t$，用 $m$、$M$ 分别表示卫星和地球的质量，$r$ 表示卫星到地心的距离，有
$G\frac{mM}{r^{2}}=mr\left(\frac{2\pi}{T}\right)^{2}$ ①

春分时，如图所示，圆 E 表示轨道，S 表示卫星，A 表示观察者，O 表示地心，由图可以看出当卫星 S 绕地心 O 转到图示位置以后（设地球自转是沿图中逆时针方向），其正下方的观察者将看不见它，据此再考虑对称性，有
$r\sin\theta=R$ ②

$t=\frac{2\theta}{2\pi}T$ ③

$G\frac{M}{R^{2}}=g$ ④

由以上各式解得 $t=\frac{T}{\pi}\arcsin\left(\frac{4\pi^{2}R}{gT^{2}}\right)^{\frac{1}{3}}$ ⑤

①式 2 分，②式 5 分，③式 5 分，④式 1 分，⑤式 3 分。

\onepicture[w=6.2cm]{figs/16_an.jpg}
}

\item
%% number: 17
%% paperName: 2004年广东广西高考第 17 题 16 分
%% typeId: 6
%% type: 计算题
%% chapter: 运动和力
%% point: 动量守恒定律
%% method: 无
%% score: 16
%% degree: 450
%% duplicateId: 0
%% body:
图中，轻弹簧的一端固定，另一端与滑块 B 相连，B 静止在水平导轨上，弹簧处在原长状态。另一质量与 B 相同滑块 A，从导轨上的 P 点以某一初速度向 B 滑行，当 A 滑过距离 $l_{1}$ 时，与 B 相碰，碰撞时间极短，碰后 A、B 紧贴在一起运动，但互不粘连。已知最后 A 恰好返回出发点 P 并停止。滑块 A 和 B 与导轨的滑动摩擦因数都为 $\mu$，运动过程中弹簧最大形变量为 $l_{2}$，求 A 从 P 出发时的初速度 $v_{0}$。

\onepicture[align=right, w=6.6cm]{figs/17.png}
%% answer: 
\jdanswer{
$v_{0}=\sqrt{\mu g(10l_{1}+16l_{2})}$
}
%% memo:
\memoanswer{
令 A、B 质量皆为 $m$，A 刚接触 B 时的速度为 $v_{1}$（碰前），由功能关系有
$\frac{1}{2}mv_{0}^{2}-\frac{1}{2}mv_{1}^{2}=\mu mgl_{1}$ ①

A、B 碰撞过程中动量守恒，令碰后 A、B 共同运动的速度为 $v_{2}$，有
$mv_{1}=2mv_{2}$ ②

碰后，A、B 先一起向左运动，接着 A、B 一起被弹回，当弹簧恢复到原长时，设 A、B 的共同速度为 $v_{3}$，在这过程中，弹簧势能始末两态都相等，利用功能关系，有
$\frac{1}{2}(2m)v_{2}^{2}-\frac{1}{2}(2m)v_{3}^{2}=\mu(2m)g(2l_{2})$ ③

此后 A、B 开始分离，A 单独向右滑到 P 点停下，由功能关系有
$\frac{1}{2}mv_{3}^{2}=\mu mgl_{1}$ ④

由以上各式，解得 $v_{0}=\sqrt{\mu g(10l_{1}+16l_{2})}$ ⑤

①②式各 2 分，③式 6 分，④⑤式各 3 分。
}

\item
%% number: 18
%% paperName: 2004年广东广西高考第 18 题 17 分
%% typeId: 6
%% type: 计算题
%% chapter: 磁场
%% point: 洛伦兹力
%% method: 无
%% score: 17
%% degree: 450
%% duplicateId: 0
%% body:
如图，真空室内存在匀强磁场，磁场方向垂直于纸面向里，磁感应强度的大小 $B=0.60\UT$，磁场内有一块平面感光板 ab，板面与磁场方向平行，在距 ab 的距离 $l=16\Ucm$ 处，有一个点状的 $\alpha$ 放射源 S，它向各个方向发射 $\alpha$ 粒子，$\alpha$ 粒子的速度都是 $v=3.0\times10^{6}\Ums$，已知 $\alpha$ 粒子的电荷与质量之比 $\frac{q}{m}=5.0\times10^{7}\UCkg$，现只考虑在图纸平面中运动的 $\alpha$ 粒子，求 ab 上被 $\alpha$ 粒子打中的区域的长度。

\onepicture[align=right, w=5.8cm]{figs/18.png}
%% answer: 
\jdanswer{
$20\Ucm$
}
%% memo:
\memoanswer{
$\alpha$ 粒子带正电，故在磁场中沿逆时针方向做匀速圆周运动，用 $R$ 表示轨道半径，有
$qvB=m\frac{v^{2}}{R}$ ①

由此得 $R=\frac{v}{\left(\frac{q}{m}\right)B}$

代入数值得 $R=10\Ucm$，可见 $2R>l>R$。

因朝不同方向发射的 $\alpha$ 粒子的圆轨迹都过 S，由此可知，某一圆轨迹在图中 N 左侧与 ab 相切，则此切点就是 $\alpha$ 粒子能打中的左侧最远点，为定出 $P_{1}$ 点的位置，可作平行于 ab 的直线 cd，cd 到 ab 的距离为 $R$，以 S 为圆心，$R$ 为半径，作弧交 cd 于 Q 点，过 Q 点作 ab 的垂线，它与 ab 的交点即为 $P_{1}$，由图中几何关系得
$NP_{1}=\sqrt{R^{2}-(l-R)^{2}}$ ②

再考虑 N 的右侧，任何 $\alpha$ 粒子在运动过程中离 S 的距离不可能超过 $2R$，以 $2R$ 为半径，S 为圆心作圆，交于 N 右侧的 $P_{2}$ 点，此即右侧能打到的最远点。由图中几何关系得
$NP_{2}=\sqrt{(2R)^{2}-l^{2}}$ ③

所求长度为
$P_{1}P_{2}=NP_{1}+NP_{2}$ ④

代入数值
$P_{1}P_{2}=20\Ucm$ ⑤

①式 3 分，②③式各 5 分，④式 1 分，⑤式 3 分。

\onepicture[w=6.8cm]{figs/18_an.jpg}
}

\end{enumerate}

\end{document}
